% ===========================================================================
%  Introducción
% ===========================================================================
\chapter*{Introducci\'on}
\addcontentsline{toc}{chapter}{Introducci\'on}
\markboth{Introducci\'on}{Introducci\'on}
\label{cap:intro}

% ---------------------------------------------------------------------------
\section*{Contexto}
% ---------------------------------------------------------------------------

La gesti\'on de la informaci\'on acad\'emica en una facultad universitaria
produce, de manera recurrente, un tipo particular de documento: la hoja de
c\'alculo de planificaci\'on.  Los horarios de docencia, los calendarios de
defensas de tesis, los cuadros de distribuci\'on de aulas y las
programaciones de actividades extracurriculares comparten una estructura
com\'un: tablas relacionadas, con columnas que resumen informaci\'on de
otras tablas mediante f\'ormulas, y con celdas que se colorean
autom\'aticamente cuando se detecta un conflicto o una condici\'on
relevante~\cite{ldmag2025}.

La hoja de c\'alculo es la herramienta elegida para estos documentos por
razones pr\'acticas: est\'a disponible en cualquier puesto de trabajo, no
requiere instalaci\'on de software adicional y produce un resultado que el
usuario final puede editar directamente.  Sin embargo, cuando el documento
debe generarse autom\'aticamente a partir de datos externos ---una base de
datos, un formulario, un archivo JSON--- la simplicidad operativa de la hoja
de c\'alculo se convierte en complejidad de programaci\'on.

% ---------------------------------------------------------------------------
\section*{El problema}
% ---------------------------------------------------------------------------

El problema central que este trabajo aborda es la mezcla de responsabilidades
en el c\'odigo que genera hojas de c\'alculo complejas.

Para ilustrarlo, consideremos un caso concreto: construir autom\'aticamente
un horario de defensas de tesis en Excel.  El documento debe incluir una
tabla de defensas con las columnas de rol del tribunal, una tabla de
profesores con el n\'umero de defensas en que participa cada uno, y una regla
de formato condicional que coloree en rojo a los profesores que aparecen en
dos defensas programadas en el mismo slot horario.

Un programa imperativo que genera este documento en Python usando la
biblioteca \texttt{openpyxl} debe conocer simult\'aneamente:
\begin{itemize}
  \item que la columna \texttt{tutor} ocupa la letra B y \texttt{secretario}
    ocupa la letra F en la hoja de c\'alculo concreta;
  \item que los datos de defensas empiezan en la fila 4 y terminan en la fila 6;
  \item que la f\'ormula de conteo es
    \texttt{COUNTIF(\$B\$4:\$F\$6,K4)};
  \item que la regla de coloreado condicional es una \texttt{FormulaRule}
    con la expresi\'on
    \texttt{SUMPRODUCT((\$B\$4:\$B\$6=\$K4)+\ldots) * COUNTIFS(\ldots) > 0 > 0}.
\end{itemize}

Si el n\'umero de defensas cambia, o si se agrega una columna de rol, o si
las tablas se reorganizan, \emph{todos} esos valores deben recalcularse
manualmente.  El c\'odigo no expresa la intenci\'on ---``colorear a los
profesores con conflicto de horario''--- sino el mecanismo concreto de su
implementaci\'on en una versi\'on espec\'{\i}fica de la hoja.

Esta situaci\'on no es exclusiva del horario de defensas.  El horario docente
de la facultad, que distribuye grupos, turnos y aulas en m\'ultiples hojas
con referencias cruzadas, y el horario de programaci\'on televisiva, que
calcula el fin de cada programa a partir de su hora de inicio y duraci\'on,
presentan el mismo patr\'on: la l\'ogica del problema queda ahogada entre
detalles de la API.

La pregunta de investigaci\'on central de este trabajo es:

\medskip
\begin{center}
\itshape
?`Puede un lenguaje de dominio espec\'{\i}fico separar la intenci\'on
sem\'antica del mecanismo de implementaci\'on, permitiendo generar
hojas de c\'alculo complejas de planificaci\'on con un c\'odigo declarativo
de nivel de abstracci\'on superior?
\end{center}
\medskip

% ---------------------------------------------------------------------------
\section*{Hip\'otesis de la investigaci\'on}
% ---------------------------------------------------------------------------

La utilizaci\'on de un lenguaje de dominio espec\'{\i}fico (DSL) interno,
implementado mediante macros de Common Lisp y un compilador basado en el
Sistema de Objetos de Common Lisp (CLOS), permite separar la descripci\'on
sem\'antica de una hoja de c\'alculo de planificaci\'on ---su estructura
tabular, sus relaciones, sus f\'ormulas y sus condiciones de coloreado--- de
los detalles concretos de la API de generaci\'on de archivos \xlsx{}.  Esta
separaci\'on reduce el c\'odigo necesario, facilita la reutilizaci\'on ante
cambios en los datos y hace que el programa resultante sea legible por el
experto del dominio sin conocimiento de la API de destino.

% ---------------------------------------------------------------------------
\section*{Objetivo general}
% ---------------------------------------------------------------------------

Dise\~nar e implementar un lenguaje de dominio espec\'{\i}fico en Common Lisp
para la declaraci\'on y generaci\'on autom\'atica de hojas de c\'alculo de
planificaci\'on, incluyendo tablas relacionadas, columnas calculadas,
cuantificaci\'on existencial y formato condicional persistente en el archivo
\xlsx{}.

% ---------------------------------------------------------------------------
\section*{Objetivos espec\'{\i}ficos}
% ---------------------------------------------------------------------------

\begin{enumerate}
  \item Analizar las caracter\'{\i}sticas recurrentes de las hojas de c\'alculo
    de planificaci\'on en el contexto acad\'emico, identificando los patrones
    comunes de estructura, c\'alculo y visualizaci\'on condicional.

  \item Dise\~nar la sintaxis y la sem\'antica de un DSL declarativo capaz de
    expresar tablas relacionadas, referencias cruzadas entre tablas,
    columnas computadas, cuantificaci\'on existencial y posicionamiento
    relativo de expresiones auxiliares.

  \item Definir un \'arbol de sintaxis abstracta (AST) cuyas clases
    representen los conceptos del dominio de manera independiente del
    backend de generaci\'on.

  \item Implementar el DSL como un conjunto de macros de Common Lisp que
    construyen nodos del AST en tiempo de expansi\'on de macros.

  \item Implementar un backend de generaci\'on de c\'odigo basado en una
    funci\'on gen\'erica CLOS que traduzca cada tipo de nodo AST a
    instrucciones Python para la biblioteca openpyxl.

  \item Validar el lenguaje sobre tres casos de uso reales de la Facultad
    de Matem\'atica y Computaci\'on de la Universidad de La Habana: el
    horario de defensas de tesis, el horario docente de la facultad y
    el horario de programaci\'on televisiva.
\end{enumerate}

% ---------------------------------------------------------------------------
\section*{Estructura del documento}
% ---------------------------------------------------------------------------

El documento se organiza de la siguiente manera.

El cap\'{\i}tulo~\ref{cap:teoria} sienta las bases te\'oricas y conceptuales del
trabajo.  Se presentan las definiciones fundamentales sobre hojas de c\'alculo,
se introduce el concepto de lenguaje de dominio espec\'{\i}fico y su uso en la
construcci\'on de compiladores, se explica el papel de los \'arboles de sintaxis
abstracta como representaci\'on intermedia, y se describen las caracter\'{\i}sticas
de Common Lisp ---macros, CLOS y despacho m\'ultiple--- que hacen de \'el una
plataforma id\'onea para el dise\~no del DSL.

El cap\'{\i}tulo~\ref{cap:dominio} analiza el dominio de la planificaci\'on
tabular a partir de tres casos reales.  Se identifican los elementos comunes
en estos documentos y se formulan los requisitos concretos que el lenguaje debe
satisfacer.

El cap\'{\i}tulo~\ref{cap:dsl} describe el lenguaje de dominio espec\'{\i}fico
dise\~nado.  Se presentan el \'arbol de sintaxis abstracta con sus cuatro
familias de nodos, las formas del DSL visible al programador y su sem\'antica,
y el flujo completo desde la declaraci\'on hasta la producci\'on del
archivo \xlsx{}.

El cap\'{\i}tulo~\ref{cap:generacion} detalla el proceso de generaci\'on de
c\'odigo.  Se explica c\'omo el backend resuelve las coordenadas Excel, c\'omo
compila cada tipo de expresi\'on a cadenas de f\'ormula, y c\'omo genera las
f\'ormulas SUMPRODUCT que implementan el formato condicional persistente.

El cap\'{\i}tulo~\ref{cap:implementacion} describe los detalles de implementaci\'on:
la organizaci\'on en cuatro archivos con responsabilidades separadas, la macro
\dsl{defclass*} que automatiza la definici\'on de nodos del AST, y los
principios de extensibilidad del sistema.

El documento concluye con las conclusiones del trabajo, las recomendaciones
para investigaciones futuras y las referencias bibliogr\'aficas consultadas.
